\section{Conclusion}

Forecast accuracy remains central to demand planning, but it is not sufficient for deployment. Operational systems execute planning signals, not forecast-error tables. A model that improves WAPE can still create unstable targets, frequent adjustments, switching burden, or execution violations that make the plan less deployable.

This paper formulates forecast-driven demand planning as feasibility-constrained model selection. The framework evaluates candidate strategies using inventory exposure, planning volatility, switching burden, and execution-capacity violations alongside forecast accuracy. DataCo provides an empirical anchor for execution-risk sensitivity, Favorita provides the main forecast-to-plan proof of concept, and M5 provides hierarchical robustness evidence.

The current results show that execution-risk-sensitive objectives can reorder deployable strategy rankings. They also show that the preferred deployable method depends on execution scenario, planning granularity, and demand intermittency. This does not imply that any single feasibility-aware method universally dominates. It supports the narrower and more operationally useful conclusion that feasibility must be evaluated explicitly when forecasting models are selected for planning systems.

The manuscript is intentionally structured as a living draft. Additional Walmart and Rossmann robustness experiments, richer execution-cost calibration, and enterprise workflow evidence can be incorporated without changing the core argument: accurate forecasts become operationally valuable only when they can be converted into feasible plans.
